\section{Understanding Quadric Surfaces Through Sections}

A three-dimensional surface should not be recognised only from a memorised
formula. We study it by intersecting it with simple planes. Fixing $z=k$ gives a
horizontal section; fixing $x=k$ or $y=k$ gives a vertical section. The remaining
equation describes a familiar planar curve.

For the ellipsoid
\[
\frac{x^2}{a^2}+\frac{y^2}{b^2}+\frac{z^2}{c^2}=1,
\]
the section $z=k$ satisfies
\[
\frac{x^2}{a^2}+\frac{y^2}{b^2}=1-\frac{k^2}{c^2}.
\]
Thus $|k|<c$ gives an ellipse, $|k|=c$ gives one point, and $|k|>c$ gives no
real points. The surface is understood as ellipses that grow from one endpoint,
reach their largest size at $z=0$, and shrink to the other endpoint.

For the elliptic paraboloid
\[
z=\frac{x^2}{a^2}+\frac{y^2}{b^2},
\]
the section $z=k$ is
\[
\frac{x^2}{a^2}+\frac{y^2}{b^2}=k.
\]
There are no points for $k<0$, one vertex for $k=0$, and growing ellipses for
$k>0$. Setting $y=0$ instead gives the parabola
\[
z=\frac{x^2}{a^2}.
\]

For the hyperboloid of one sheet
\[
\frac{x^2}{a^2}+\frac{y^2}{b^2}-\frac{z^2}{c^2}=1,
\]
the section $z=k$ becomes
\[
\frac{x^2}{a^2}+\frac{y^2}{b^2}=1+\frac{k^2}{c^2}.
\]
Every height gives an ellipse, smallest at $k=0$ and larger as $|k|$ grows.
This explains the connected surface and its central waist.

\begin{remarkbox}
For every quadric below, ask which planar curve appears after one coordinate is
fixed, for which values that curve exists, and how it changes as the cutting
plane moves. The standard-form atlas records the surfaces; their sections
explain them.
\end{remarkbox}
